\scanpage{63}
\uline{Compléments}
\textcircled{1}

\uline{Pb} Pour $R\in\operatorname{Ob}\Delta$ et $\underset{\text{cat. discr.}}I\longrightarrow R$, caractériser la situation où cela correspond à une famille d'objets de base (indexée par $I$) pour une théorie \uline{algébrique} $T_R$. (Il y a une caractérisation constructive transfinie : $R$ est l'aboutissement d'une suite transfinie de catégories $R_\alpha$ sur $R$, faisant un \inserted{(pseudo)} syst. inductif \deleted{de catégories \unsure{$C_\Delta$}} dans $\Delta$, avec
\[
 \begin{aligned}
 (1)\quad &R_0=I\\
 (2)\quad &R_{\alpha+1}\ \text{déduit de $R_\alpha$ soit en rendant inversibles un petit ens. de flèches}\\
 &\hspace{7em}\text{\deleted{soit en \unsure{adjoignant|égalisant}} \inserted{de couples de flèches}}\\
 &\hspace{7em}\text{soit en ajoutant un petit ens de flèches}\\
 (3)\quad &R_\alpha=\varinjlim_{\beta<\alpha} R_\beta\qquad\text{si $\alpha$ ordinal limite.}
 \end{aligned}
\]
\margin{Condition nécessaire (?) : toute sous-catégorie pleine $\in\Delta$ de $R$ qui contient $I$, est \unsure{ess.t|ess.l.} égale à $R$. Est-ce suffisant ? \unsure{Prob.t oui|Prob.t que}, \deleted{\unsure{une}} qu'une ``\uline{petite}'' construction transfinie de $R$ à partir de $I$\ldots}% slanted
\margin{\unsure{marchent} pour $I$ c.d. petit implique $\widetilde I$ petite, et $C\in\Delta$ petite implique $C\Sigma^{-1}$ petite}% slanted

\uline{NB 1} \deleted{la construction transfinie et} la considération des $R$ non petits se pose seulement si $(d'_i,d''_i)$ n'est pas ``petit'', ou du moins de ramener à un ens petit \unsure{!\ldots!|$\mathcal E$\ldots!} Néanmoins, même si $(d'_i,d''_i)$ petits, la question de caractériser les $I\mathbin{-}R$, à partir d'une construction transfinie, se pose.

Dans la définition récurrente des théories algébriques, \deleted{\unsure{$(d$|(cas d)} \inserted{$\operatorname{Card}d_i$}} , \deleted{\unsure{NB}} il faut admettre des $\varprojlim_{\alpha}$ sur des ens. bien ordonnés\ldots.

Il faudrait néanmoins prouver que si on a des $I\mathbin{-}(R_\lambda)_{\lambda\in\Lambda}$ satisfaisant la condition envisagée (disons ``$I$ $\Delta$-engendre $R_\lambda$) alors c'est encore pareil pour $I\mathbin{-}R$, où $R$ est la 2-somme amalgamée des $R_\lambda$ sous $I$ dans $\Delta$.

\uline{NB 2} Soit $J$ la (petite) sous-catégorie pleine de $R$ engendr. par les objets $\in I$ et les \inserted{(images dans $R$ des)} sources et buts de toutes les flèches rajoutées au cours des \deleted{\unsure{mod}} dans la construction transfinie. Il est plausible que $\forall C\in\operatorname{Ob}\Delta$, $\operatorname{Hom}_\Delta(R,C)\overset{=T_R(C)}{\longrightarrow}\underline{\operatorname{Hom}}_{\mathrm{Cat}}(J,C)\simeq\underline{\operatorname{Hom}}_\Delta(\widetilde J,C)=T_{\widetilde J}(C)$ est \uline{pleinement fidèle}. \deleted{\unsure{En}} Or $T_{\widetilde J}$ est une théorie algébrique. Il est plausible \inserted{dès lors} que \deleted{\unsure{en fait}} $T_R$ se déduit de $T_{\widetilde J}$ \deleted{\unsure{par}} en rajoutant seulement des axiomes. Mais il ne suffit peut-être pas de le faire en une
% archive sheet number: 62
